\documentclass[11pt]{article}
\usepackage[margin=1in]{geometry}
\usepackage{amsmath,amssymb,amsthm,booktabs}
\usepackage[hidelinks]{hyperref}
\newtheorem{theorem}{Theorem}
\newtheorem{lemma}{Lemma}
\newcommand{\asinh}{\operatorname{arsinh}}
\newcommand{\R}{\mathcal R}
\title{The strict \(M_4\) resultant-mean inequality}
\author{Principia Math In Progress}
\date{18 September 2026}
\begin{document}
\maketitle
\begin{abstract}
We prove both parameter branches of Mihokovi\'c--Ra\"issouli's
\(M_4\) sign conjecture. A mixed derivative reduces the assertion to
positivity of one explicit function. Analytic estimates cover a
neighborhood of zero and the entire infinite tail; a reproducible
interval-arithmetic certificate covers the intervening compact interval.
The finite computation encloses whole intervals, rather than sampling
points.
\end{abstract}

\section*{Task 1: the original problem}
For \(a,b>0\), define
\[
 M(a,b)=M_4(a,b)=
 \frac{b-a}{\sqrt2\,\asinh((\log b-\log a)/\sqrt2)}
 \quad(a\ne b),\qquad M(a,a)=a.
\]
For \(r\ne0\), let \(B_r(a,b)=((a^r+b^r)/2)^{1/r}\), and put
\[
 \R(B_p,M,B_q)(a,b)
 =B_p\bigl(M(a,B_q(a,b)),M(B_q(a,b),b)\bigr).
\]
The intended strict comparison is for distinct positive inputs; the
diagonal has equality. This is Conjecture~4.10 of the preprint
\cite{MR}, and Conjecture~40, p.~22, of its complete published version.

\begin{theorem}\label{main}
If \(p=\sqrt{21}\) or \(p=-\sqrt{21}\), and \(q=(3-p)/2\), then
\[
 M(a,b)>\R(B_p,M,B_q)(a,b)\qquad(a,b>0,\ a\ne b).
\]
Thus both choices \(q=(3\pm\sqrt{21})/2\), \(p=3-2q\), in the
conjecture are valid.
\end{theorem}

\section*{A mixed-derivative reduction}
By symmetry, assume \(a<b\), and set
\[
 H(a,b)=(b^q-a^q)M(a,b)^p,\qquad H(a,a)=0.
\]
We will prove \(H_{ab}>0\) for \(a<b\).
For any \(a<c<b\), the fundamental theorem of calculus then gives
\begin{equation}\label{rectangle}
 H(a,c)+H(c,b)-H(a,b)
   =\int_a^c\int_c^b H_{ab}(u,v)\,dv\,du>0.
\end{equation}
Take \(c=B_q(a,b)\), so that
\(c^q-a^q=b^q-c^q=(b^q-a^q)/2\).
If \(q>0\), \eqref{rectangle} gives
\(M(a,b)^p<(M(a,c)^p+M(c,b)^p)/2\); here \(p<0\), giving the
claimed strict comparison. If \(q<0\), the first inequality reverses
on division by \(b^q-a^q\); now \(p>0\), giving the same conclusion.

Here is an explicit test for the mixed derivative. For \(t>0\), write
\[
 v(t)=\log\frac{\sqrt2\sinh(t/2)}{\asinh(t/\sqrt2)},\quad
 g(t)=\frac1{\sqrt{t^2+2}\,\asinh(t/\sqrt2)},\quad s=p+q.
\]
Then
\begin{align}
 v'(t)&=\tfrac12\coth(t/2)-g(t),\label{vprime}\\
 v''(t)&=-\frac1{4\sinh^2(t/2)}
              +\frac{t}{t^2+2}g(t)+g(t)^2.\label{vsecond}
\end{align}
Define
\begin{equation}\label{T}
 T_p(t)=v''(t)+p\,v'(t)^2+q\coth(qt/2)v'(t)-\frac34.
\end{equation}
With \(t=\log(b/a)\), we have
\[
 H(a,b)=a^s h(t),\qquad
 h(t)=e^{st/2}\,2\sinh(qt/2)e^{pv(t)}.
\]
Since \(p+2q=3\), direct differentiation gives
\[
 h''-s h'=p h T_p,\qquad
 H_{ab}(a,b)=-\frac{a^s}{ab}\,p h(t)T_p(t).
\]
Indeed, the constant in the first identity is
\((q^2-s^2)/4=-3p/4\).
The sign of \(h(t)\) is the sign of \(q\), and \(pq<0\).
It remains to prove \(T_p(t)>0\) for \(t>0\), for both values of \(p\).

\section*{A rigorous certificate for \(T_p>0\)}
All interval computations below use outward enclosures in Arb, through
python-flint \cite{arb}, at 160-bit midpoint precision. The accompanying
\texttt{verify\_inequality.py} recomputes the coefficients, partitions,
and strict enclosures, and writes \texttt{exact-certificate.json}.
The analytic estimates below supply the two infinite families of
arguments outside the compact computation.

\subsection*{1. An analytic Taylor remainder at zero}
The function \(v'\) is odd and analytic near zero. Its expansion begins
\[
 v'(t)=\frac t4-\frac{t^3}{16}+\frac{17t^5}{672}+O(t^7).
\]
Consequently \(T_p\) is even and analytic, and, using \(p^2=21\),
\begin{equation}\label{leading}
 T_p(t)=\frac{13}{480}t^4+O(t^6).
\end{equation}
For explicit verification, the coefficients of \(t^2\) and \(t^4\) are
\[
 \frac{-15+3p+2q^2}{48}=0,\qquad
 \frac{255-15q^2-q^4-45p}{1440}=\frac{13}{480}.
\]
The constant coefficient is zero as well.

We need a uniform, proved bound on all omitted coefficients. Define
\[
 \rho(z)=\frac{\asinh(z/\sqrt2)}{z/\sqrt2},\qquad
 U(z)=\frac1{\sqrt{1+z^2/2}\,\rho(z)},\qquad
 D(w)=w\coth w.
\]
The removable values at zero are one. For \(|z|\le1\), the binomial
series gives
\[
 |\rho(z)-1|\le\frac13,\qquad
 |(1+z^2/2)^{-1/2}-1|\le\sqrt2-1<\frac12.
\]
Thus \(|U-1|<5/4\). Also, on \(|z|\le1\),
\[
 |D(z/2)-1|
 \le\frac{(\cosh(1/2)-1)+(\sinh(1/2)/(1/2)-1)}
 {1-(\sinh(1/2)/(1/2)-1)}<\frac14.
\]
For example, the two numerator terms are less than \(1/7\) and \(1/20\).
The denominators in these formulas have no zeros in the specified
discs. Since \(v'(z)=(D(z/2)-U(z))/z\), the maximum principle gives
\[
 |v'(z)/z|<\frac32\quad(|z|\le1).
\]
Cauchy's derivative estimate therefore implies, on \(|z|\le1/2\),
\[
 |v'(z)|<\frac34,\qquad |v''(z)|<3.
\]
Furthermore, \(|q|<4\), and for \(|w|\le1\),
\[
 |D(w)|\le\frac{\cosh1}{1-(\sinh1-1)}<2.
\]
Here \(\cosh1<8/5\) and \(\sinh1-1<1/5\).
Writing \(q\coth(qz/2)v'(z)=2D(qz/2)v'(z)/z\), we obtain
\[
 |T_p(z)|<3+5(3/4)^2+6+3/4<16\qquad(|z|\le1/2).
\]
These estimates also establish analyticity throughout that disc.
If \(T_p(t)=\sum a_k t^k\), Cauchy's bound is
\(|a_k|\le16\,2^k\).
Parity and \eqref{leading} now give, for \(0<t\le1/4\),
\begin{equation}\label{tailzero}
 \left|\frac{T_p(t)-\sum_{k=2}^{20}a_{2k}t^{2k}}{t^4}\right|
 \le\frac{16\,2^{42}t^{38}}{1-2t}\le 2^{-29}.
\end{equation}

The verifier computes the coefficients through degree 40 by formal
ball-series arithmetic on the displayed formulas. Horner evaluation
of \(\sum_{k=2}^{20}a_{2k}y^{k-2}\) on the single interval
\(0\le y\le1/16\), followed by subtraction of \(2^{-29}\), gives a
lower bound greater than \(1/40\), for each branch. It follows that
\[
 T_p(t)>\frac{t^4}{40}>0\qquad(0<t\le1/4).
\]
Formal-series division by \(t\) removes only the exact zero terms
justified by parity and the displayed removable values.

\subsection*{2. The compact interval}
Start with the exact rational interval \([1/4,20]\). For each interval
\([l,u]\), evaluate \eqref{vprime}--\eqref{T} on the Arb ball centered
at \((l+u)/2\) with radius \((u-l)/2\). Accept the interval only when
the entire resulting enclosure is strictly positive; otherwise bisect
it at its exact rational midpoint and repeat. The computation terminates:
\[
\begin{array}{c|r|c}
 p&\text{accepted subintervals}&\text{certified lower bound on every leaf}\\
 \hline
 -\sqrt{21}&186338&10^{-12}\\
 \sqrt{21}&173452&10^{-12}
\end{array}
\]
The code checks that the exact sum of accepted lengths is \(79/4\).
The binary subdivision preserves a cover with disjoint interiors;
thus this is a certificate for every point of \([1/4,20]\).
The JSON certificate records partition hashes and depth counts.
No unverified decimal approximation of \(p\), \(q\), or an endpoint
is used.

\subsection*{3. The whole tail \(t\ge20\)}
Let \(r=\sqrt{21}\), \(k=(r-3)/2\), and
\(\epsilon=(e^t-1)^{-1}\). We have
\[
 \frac{3}{4}<k<1,\quad 4<r<5,\quad
 v'=\frac12-g+\epsilon,\quad
 v''=g^2+\frac{t}{t^2+2}g-\epsilon(1+\epsilon).
\]
For \(t\ge20\),
\[
 \frac1{2t\log(2t)}<g<\frac1t\le\frac1{20},
 \qquad 0<\epsilon<2e^{-t}.
\]
The first lower bound follows from
\(\sqrt{t^2+2}<2t\) and \(\asinh(t/\sqrt2)<\log(2t)\).
In particular \(v'>0\).

For \(p=-r\), \(q=(3+r)/2\), use \(\coth(qt/2)>1\) in \eqref{T}
and expand. Dropping positive terms gives
\begin{align*}
 T_{-r}
 &\ge kg-(r-1)g^2-(k+1)\epsilon-(r+1)\epsilon^2\\
 &>\frac12g-2\epsilon-6\epsilon^2\\
 &>\frac1{4t\log(2t)}-4e^{-t}-24e^{-2t}\\
 &\ge\frac1{4t^2}-28e^{-t}>0.
\end{align*}
Here \(k-(r-1)g>3/4-4/20>1/2\), and
\(\log(2t)<t\). For the final strict inequality,
\(t^2e^{-t}\) decreases for \(t\ge20\), while
\[
 e^{20}>\frac{20^8}{8!}>44800=112\cdot20^2.
\]

For \(p=r\), \(q=-k\), we instead have
\(q\coth(qt/2)>k\) and \(v'>9/20\). Consequently
\[
 T_r>-\epsilon(1+\epsilon)
       +4(9/20)^2+\frac34(9/20)-\frac34
 =-\epsilon(1+\epsilon)+\frac{159}{400}>0.
\]
Indeed \(\epsilon<2e^{-20}<1/100\), so its negative contribution
is less than \(101/10000\).
This completes the proof that \(T_p>0\) for every \(t>0\).
The mixed-derivative reduction and \eqref{rectangle} prove
Theorem~\ref{main}. \(\square\)

\paragraph{Scope of the proof and audit.}
This is a computer-assisted proof: its compact and near-zero
polynomial enclosures are substantive proof components.
The script is independently rerunnable and supplies rigorous
enclosures, not merely high-precision samples.
The authoring agent checked the algebra, analytic tails, source scope,
and executed the certificate. Independent peer review and
proof-assistant formalization are not claimed.

\section*{Proposed follow-up questions}
The following four questions are unproved directions, not additional solves.
\paragraph{Task 2.}
Find an elementary analytic proof of \(T_{\pm\sqrt{21}}(t)>0\)
that removes the interval computation.
\paragraph{Task 3.}
Determine the complete region of real parameter pairs \((p,q)\) for
which the inequality
\(M_4>\R(B_p,M_4,B_q)\) holds for every distinct positive pair.
\paragraph{Task 4.}
Prove or disprove the analogous \(M_2\) inequality in Conjecture~37
of the published source, for both prescribed power-mean pairs.
\paragraph{Task 5.}
For each branch here, determine the largest constant \(C_p>0\) such that
\[
 \frac{M_4(a,b)-\R(B_p,M_4,B_q)(a,b)}{M_4(a,b)}
 \ge C_p\left(\frac{|\log(b/a)|}{1+|\log(b/a)|}\right)^6
\]
for every distinct positive pair, and determine whether the infimum
is approached on the diagonal.

\begin{thebibliography}{9}
\raggedright
\bibitem{MR}
L.~Mihokovi\'c and M.~Ra\"issouli,
\emph{Results on comparison and sub/super-stabilizability of some new means},
Applicable Analysis and Discrete Mathematics \textbf{20} (2026), 1--23,
\href{https://doi.org/10.2298/AADM240618001M}{doi:10.2298/AADM240618001M}.
Complete published text:
\href{http://pefmath.etf.rs/vol20num1/AADM-Vol20-No1-1-23.pdf}{journal PDF}.
Preprint: \href{https://arxiv.org/abs/2405.13618}{arXiv:2405.13618v2}.
\bibitem{arb}
The python-flint developers,
\emph{Arb real numbers and rigorous ball arithmetic},
\href{https://python-flint.readthedocs.io/en/latest/arb.html}{official documentation}.
The accompanying certificate records the exact installed version.
\end{thebibliography}
\end{document}
